\section{Introduction}
\label{sec:introduction}

What does a student know, and what will they do next?  Educational systems
usually operationalize this question as a collection of separate prediction
problems.  Knowledge tracing estimates whether the next response will be
correct \citep{corbett1995knowledge,piech2015deep,ghosh2020context}; cognitive
diagnosis associates performance with latent skills; response models predict
a selected option; and misconception systems explain a particular error.
These targets are related, but the models used to predict them rarely share an
interface.  A system that outputs a scalar probability of correctness cannot
directly say that a student is likely to add two denominators, produce the
resulting fraction, or preserve that error on a later problem without adding a
new task-specific output space.

This separation matters because correctness is a lossy description of student
behavior.  Two students can both answer incorrectly while producing different
responses for different reasons, and those differences determine which future
errors are plausible and which feedback would be appropriate.  Recent datasets
make richer modeling possible: Eedi links distractors to expert-authored
misconceptions \citep{wang2021diagnostic}; MalruleLib provides executable
mathematical error rules \citep{chen2026malrulelib}; and FoundationalASSIST
contains longitudinal problem text together with students' first responses
\citep{worden2026foundationalassist}.  Yet the dominant knowledge-tracing
formulation still maps an interaction sequence to a fixed correctness head.

Large language models (LLMs) offer a natural way to express heterogeneous
targets as conditional generation.  They can consume problem text, answer
choices, skill descriptions, and student histories in a shared context, then
emit either a class label or an open-form response.  However, using an LLM does
not by itself produce a valid student model.  A capable language model may
solve the problem rather than imitate the student, rely on item priors rather
than the student's history, or generate a plausible explanation unrelated to
the student's observed error.  Prior work has accordingly questioned whether
LLMs can simulate students under perturbed evidence \citep{sonkar2023deduction}
and whether apparent student simulations remain valid under stronger controls
\citep{yuan2026valid}.  A useful generative student model must therefore pass
both predictive and mechanistic tests: it should use student-specific evidence
to predict multiple outcomes, and its internal student representation should
have a demonstrable effect on those predictions.

We introduce \model, a causally grounded conditional-generation model with an
explicit recurrent student state.  After each observed interaction, a shared
language-model updater transforms the previous continuous state and the new
evidence into $z_t$.  A task-conditioned decoder then receives only $z_t$ and a
query; the raw history is unavailable as a bypass.  The same model produces
structured outputs for five constructs:
(i)~next-response correctness, (ii)~the student's concrete response,
(iii)~the misconception expressed by an observed error, (iv)~correctness and
response on a future problem after unobserved intervening interactions, and
(v)~the recurrence of the same wrong response on a later exposure.  The
formulation includes both closed-choice and open-generation protocols; it does
not average their unlike metrics into a single ``student modeling score.''

The main methodological challenge is to make $z_t$ sufficient for prediction
without sacrificing the performance of a full-context language model.  We use
four complementary objectives.  First, a full-context teacher distills its
output distribution into the state-only decoder.  Second, a trajectory
consistency loss aligns recurrently updated states with states independently
encoded from the corresponding complete history.  Third, executable
misconceptions provide counterfactual supervision: swapping a donor student's
state into a recipient query should generate the error obtained by applying the
donor's rule.  The last objective grounds state interventions behaviorally
rather than training an auxiliary misconception classifier.  Fourth, a
future-profile objective aligns a frozen view of the state with later
knowledge-component outcomes while teaching matched student--state ownership.

We evaluate the resulting state with a student-disjoint longitudinal protocol.
A readout selected only on training and validation students predicts held-out
future knowledge-component rates and emits readable weakness profiles.  We
compare it with a population prior and a same-window correctness/count
baseline, then permute state ownership across test students and compare two
disjoint views of each student's observed history.  The permutation asks
whether the right student's state matters; the disjoint-view test asks whether
the recovered information is stable.  Literal cross-student text transport is
reported as a diagnostic rather than a primary gate because an arbitrary
transplanted latent trajectory need not be in distribution.  This distinction
is essential because decodability alone does not establish a causal role
\citep{hewitt2019designing,elazar2021amnesic}.

We instantiate the framework with a 32B-parameter language model and a mixture
of six data sources spanning real longitudinal interaction logs, expert-labeled
misconceptions, and executable synthetic error rules.  Our formal evaluation
keeps students disjoint across longitudinal splits and separates temporal
prediction, non-temporal profile prediction, full-taxonomy misconception
retrieval, open generation, and controlled synthetic generalization.  This
protocol is deliberately stricter than evaluating against a row-specific set
of candidates known to contain the gold answer.  It also makes the limits of
each construct explicit: future performance within a knowledge component is
not treated as a direct mastery label, and repeated incorrectness without a
recorded response is not treated as persistence of the same misconception.

Our contributions are:
\begin{itemize}
    \item \textbf{A recurrent generative student state.} We cast correctness,
    concrete response prediction, misconception diagnosis, prospective
    behavior, and exact wrong-response recurrence as task-conditioned
    generation from one continuously updated state.
    \item \textbf{A construct-aware multi-source evaluation.} We align six
    educational data sources while preserving source-specific semantics and
    compare the unified model with classical and neural specialists, including
    IRT, PFA, BKT, DKT, AKT, retrieval, and response-frequency models.
    \item \textbf{Transfer across synthetic and authentic errors.} We test
    whether supervision from executable misconception rules improves diagnosis
    of expert-labeled errors, with task-matched ablations that distinguish rule
    supervision from merely adding synthetic responses.
    \item \textbf{Grounded state training.} Full-context distillation preserves
    predictive information, trajectory consistency teaches incremental
    updating, executable counterfactuals supervise task behavior, and future
    profile supervision binds states to student-specific outcomes.
    \item \textbf{A reusable, audited student state.} We combine held-out future
    prediction, readable profiles, state-ownership permutations, matched
    baselines, and disjoint-history stability to test whether the learned state
    is useful and student-specific.
\end{itemize}

Together, these contributions shift the question from whether a language model
can predict one educational label to whether one model can represent a
student's behavior at several levels of specificity---and whether the internal
state supporting those predictions can survive causal scrutiny.
